\section{Attack Surface 2: TIGER/Line Shapefiles}
\label{sec:tiger}

\subsection{The Data Format}

TIGER/Line shapefiles define the geographic boundaries of each census tract
as a sequence of latitude/longitude coordinate pairs.
Two tracts are adjacent if their boundary polygons share a segment of
positive length (the ``shared edge'' criterion used by \bisect{}).
The adjacency graph $G = (V, E)$ is constructed from these shapefiles:
$V$ is the set of census tracts, and $(u, v) \in E$ iff tracts $u$ and $v$
share a boundary segment.

\subsection{The Attack}

A hostile actor who can modify TIGER/Line shapefiles can rewire the
adjacency graph without changing population counts.
Specifically:
\begin{enumerate}
\item \textbf{Adjacency insertion}: Slightly modify the boundary
  coordinates of two non-adjacent tracts so that they appear to share
  a boundary segment.
  This adds an edge to the adjacency graph, potentially enabling
  a bisection cut that was previously infeasible.
\item \textbf{Adjacency deletion}: Remove boundary coordinates shared
  between two adjacent tracts, breaking their adjacency.
  This removes an edge, forcing the bisection to find alternative cuts.
\item \textbf{Boundary shift}: Move a boundary line between two tracts
  (one heavily Democratic, one heavily Republican) by a small amount,
  effectively moving some people from one tract to the other
  without changing the total population.
\end{enumerate}

Attack~3 (boundary shift) is the most potent because it operates below
the GEOID level: the adversary does not change which GEOID a tract has or
what its total population is, only where its geographic boundary lies.
If the algorithm uses centroid-based population assignment (which
\bisect{} does not---it uses census-provided population counts per GEOID),
this attack could effectively move people between tracts.

\subsection{Why Bisect Is Resistant to Boundary Shifts}

The critical design decision in \bisect{} is that population is taken from
the P.L.\ 94-171 file keyed by GEOID, not derived from shapefile geometry.
A shapefile boundary shift does not change the P.L.\ 94-171 GEOID record
and therefore does not change the population count used by the algorithm.

The only effect of a shapefile boundary shift is to change the adjacency
relationship between tracts.
This is the adjacency insertion/deletion attack described above.

\subsection{Effect Magnitude}

Adding or removing individual adjacency edges in the census tract graph
has limited partisan effect because:
\begin{itemize}
\item Congressional district boundaries are determined by sequences of
  hundreds to thousands of edges in the graph; single-edge changes
  affect only the local margin of a bisection cut.
\item The ConvergenceSweep algorithm~\citep{b16paper} optimises over
  many possible cuts; a single spurious edge is unlikely to become
  the preferred cut unless it is in a highly competitive region.
\item TIGER/Line shapefiles use high-precision coordinate pairs
  (6 decimal places of latitude/longitude, approximately 10-meter
  precision).
  Inserting a false shared edge requires introducing coordinate pairs
  that are inconsistent with the surrounding polygon geometry;
  a geometric consistency check can detect such insertions.
\end{itemize}

We estimate the maximum effect of a plausible (undetected) shapefile
manipulation at $\Delta D \leq 0.4$ seats nationally.
A manipulation producing a larger effect would require altering dozens of
adjacency relationships in competitive states, which is geometrically
implausible (non-convex polygon corners, sub-meter boundary precision).

\subsection{Detection}

\texttt{bisect} computes the adjacency graph hash:
\begin{verbatim}
SHA256(adjacency_graph_state_XX_2020.json) -> manifest
\end{verbatim}
The adjacency graph is a deterministic function of the TIGER/Line shapefile.
Any modification to the shapefile that changes adjacency relationships
changes the adjacency hash.

Detection probability: $\delta = 0.98$.
The 2\% residual corresponds to modifications that change shapefile
coordinates but do not change adjacency relationships (e.g., interior
smoothing of non-boundary coordinates), which are genuinely undetectable
but also have zero effect on the algorithm's output.
